\documentclass[11pt,a4paper]{article}
\usepackage[utf8]{inputenc}
\usepackage[T1]{fontenc}
\usepackage{lmodern}
\usepackage[margin=1in]{geometry}
\usepackage{microtype}
\usepackage{setspace}
\usepackage{hyperref}
\usepackage{amsmath}
\usepackage{booktabs}
\usepackage{graphicx}
\usepackage{xcolor}
\usepackage{enumitem}
\usepackage{listings}
\usepackage{float}

\lstset{
  basicstyle=\ttfamily\small,
  breaklines=true,
  frame=single,
  backgroundcolor=\color{gray!8},
  language=Python
}

\title{\textbf{}\\
Learning and Reasoning for Legal AI via Intelligent Multi-Model Routing with a smart Evaluation Layer}
\author{Aditya Mukhopadhyay, Sara Ghane, Marco D. Gennaro}
\date{\today}

\begin{document}
\maketitle
\onehalfspacing

\begin{abstract}
Legal AI systems must handle both high-recall extraction and high-precision legal reasoning, yet most deployments still use a single language model across all tasks. This creates a persistent quality-cost mismatch: smaller models often perform well on retrieval-heavy operations, while larger models are more reliable for interpretive legal judgments. We present a working paper on a hybrid architecture that integrates learning and reasoning through interpretable task routing. The system combines rule-based signals (intent keywords, document-length heuristics, confidence fallbacks) with data-driven routing via supervised semantic features to dynamically select between a specialist extraction model and a generalist reasoning model. Our current implementation is evaluated on tender/contract extraction and LegalBench-style reasoning tasks using metrics (F2/Jaccard for extraction; accuracy-based scores for reasoning). We report current limitations and outline next steps on robustness testing, calibration, and human-centered error analysis for trustworthy legal decision support. Our pipeline currently includes 2,010 model-level evaluations over 1,005 task instances (999 reasoning, 6 extraction) and seeded router analysis on the 1,005-task set. The single-model baselines show limitations: Maverick and 70B achieve mean scores of 0.629 and 0.571 overall, but both perform weakly on extraction (0.112 and 0.100, respectively), and both have about 14\% zero-score cases.
\end{abstract}

\end{document}
